% --- Archon physics formalization source begin ---
% archon:physics
% archon:covers PhyXMiniProblems/problem_phyx_mini_0615.lean
% archon:source-report reports/phyx_mini/problem_phyx_mini_0615.source.json

\chapter{Physics problem phyx\_mini\_0615}
\label{ch:PhyXMiniProblems_problem_phyx_mini_0615}

\paragraph{Problem source.}
\#\# Physical scenario

In a simple model for a radioactive nucleus, an alpha particle (\$m = 6.64 	imes 10\textasciicircum{}\{-27\}\$ kg) is trapped by a square barrier that has width 2.0 fm and height 30.0 MeV.The energy of the alpha particle is 10.0 MeV below the top of the barrier.

\#\# Question

What is the tunneling probability?

\#\# Answer choices

- A: A: 0.014

- B: B: 0.023

- C: C: 0.09

- D: D: 0.031

\#\# Figure caption (auxiliary; use the image as primary evidence)

The image is a plot featuring a potential energy diagram:

- **Axes**: 
  - The vertical axis is labeled \textbackslash{}(U(r)\textbackslash{}).
  - The horizontal axis is labeled \textbackslash{}(r\textbackslash{}).

- **Graph**:
  - A blue stepped line represents the potential barrier.
  - The potential starts at zero, rises vertically to a level marked \textbackslash{}(U\_0\textbackslash{}), forming a rectangular shape, and then drops back down.

- **Lines**:
  - A dashed horizontal line labeled \textbackslash{}(E\textbackslash{}) runs parallel below \textbackslash{}(U\_0\textbackslash{}).
  - The difference between \textbackslash{}(E\textbackslash{}) and \textbackslash{}(U\_0\textbackslash{}) is annotated as \textbackslash{}(1.0 \textbackslash{}, \textbackslash{}text\{MeV\}\textbackslash{}).

- **Annotations**:
  - The base of the rectangle spans \textbackslash{}(2.0 \textbackslash{}, \textbackslash{}text\{fm\}\textbackslash{}).
  - Origin labeled as \textbackslash{}(O\textbackslash{}).

The plot represents a one-dimensional potential energy barrier, commonly used in physics to illustrate concepts such as tunneling.

\#\# Dataset metadata

Quantum Phenomena; Numerical Reasoning, Physical Model Grounding Reasoning

\paragraph{Recorded answer/context.}
A: A: 0.014

\paragraph{Figure/image path.}
/root/proposal\_for\_physic/science-mango/phyx\_mini\_run/phyx\_data/test\_image/615.png

\paragraph{Formalization target.}
create a compiling Lean file with sorry bodies at `PhyXMiniProblems/problem\_phyx\_mini\_0615.lean`. The Lean declarations must preserve the physical quantities, dimensions or dimensional roles, figure labels, governing-law hypotheses, and final relation expressed by this problem.
Use Mathlib/Physlib names found through LeanExplore where available. If a physics API is missing, introduce faithful local abstractions rather than scalar placeholder aliases.

\begin{theorem}[Physics formalization target]
\label{thm:physics:phyx_mini_0615:target}
\lean{PhyXMiniProblems.ProblemPhyXMini0615.problem_phyx_mini_0615}
\uses{def:physics:phyx-mini-0615:phyxminiproblems-problemphyxmini0615-energyinmegaelectronvolts, def:physics:phyx-mini-0615:phyxminiproblems-problemphyxmini0615-alpharectangularbarriersetup, def:physics:phyx-mini-0615:phyxminiproblems-problemphyxmini0615-matchesalphanucleusscenario, def:physics:phyx-mini-0615:phyxminiproblems-problemphyxmini0615-matchessuppliedbarrierfigure, def:physics:phyx-mini-0615:phyxminiproblems-problemphyxmini0615-matchesproblemreadouts, def:physics:phyx-mini-0615:phyxminiproblems-problemphyxmini0615-usesstandardreducedplanckconstant, def:physics:phyx-mini-0615:phyxminiproblems-problemphyxmini0615-matchesphyslibrectangularbarrier, def:physics:phyx-mini-0615:phyxminiproblems-problemphyxmini0615-hasphysicaltunnelingparameters, def:physics:phyx-mini-0615:phyxminiproblems-problemphyxmini0615-obeystextbookrectangularbarriertunneling, def:physics:phyx-mini-0615:phyxminiproblems-problemphyxmini0615-roundstothreedecimalplaces, def:physics:phyx-mini-0615:phyxminiproblems-problemphyxmini0615-isuniquematchingtunnelingchoice}
The assigned autoformalize agent should translate this physics problem into Lean declarations in the covered file, with theorem and lemma proof bodies written as `by sorry`.
\end{theorem}
\begin{proof}
This is an autoformalization task, not a proof task. Produce statements that can later be proved by the physics prover without weakening the physical model.
\end{proof}

% --- Archon declaration topology begin ---
\section{Declaration topology}
\begin{definition}[Length Quantity]
  \label{def:physics:phyx-mini-0615:phyxminiproblems-problemphyxmini0615-lengthquantity}
  \lean{PhyXMiniProblems.ProblemPhyXMini0615.LengthQuantity}
  A nonnegative, unit-independent physical length
\end{definition}
\begin{proof}
  This is the declared model definition of Length Quantity.
\end{proof}
\begin{definition}[Mass Quantity]
  \label{def:physics:phyx-mini-0615:phyxminiproblems-problemphyxmini0615-massquantity}
  \lean{PhyXMiniProblems.ProblemPhyXMini0615.MassQuantity}
  A nonnegative, unit-independent physical mass
\end{definition}
\begin{proof}
  This is the declared model definition of Mass Quantity.
\end{proof}
\begin{definition}[Energy Quantity]
  \label{def:physics:phyx-mini-0615:phyxminiproblems-problemphyxmini0615-energyquantity}
  \lean{PhyXMiniProblems.ProblemPhyXMini0615.EnergyQuantity}
  A signed, unit-independent physical energy
\end{definition}
\begin{proof}
  This is the declared model definition of Energy Quantity.
\end{proof}
\begin{definition}[action Dimension]
  \label{def:physics:phyx-mini-0615:phyxminiproblems-problemphyxmini0615-actiondimension}
  \lean{PhyXMiniProblems.ProblemPhyXMini0615.actionDimension}
  The dimension of action, mass * length\textasciicircum{}2 / time
\end{definition}
\begin{proof}
  This is the declared model definition of action Dimension.
\end{proof}
\begin{definition}[Action Quantity]
  \label{def:physics:phyx-mini-0615:phyxminiproblems-problemphyxmini0615-actionquantity}
  \lean{PhyXMiniProblems.ProblemPhyXMini0615.ActionQuantity}
  \uses{def:physics:phyx-mini-0615:phyxminiproblems-problemphyxmini0615-actiondimension}
  A nonnegative physical action, used for the reduced Planck constant
\end{definition}
\begin{proof}
  This is the declared model definition of Action Quantity.
\end{proof}
\begin{definition}[Wave Number Quantity]
  \label{def:physics:phyx-mini-0615:phyxminiproblems-problemphyxmini0615-wavenumberquantity}
  \lean{PhyXMiniProblems.ProblemPhyXMini0615.WaveNumberQuantity}
  A nonnegative physical wave number, carrying inverse-length dimension
\end{definition}
\begin{proof}
  This is the declared model definition of Wave Number Quantity.
\end{proof}
\begin{definition}[length Readout]
  \label{def:physics:phyx-mini-0615:phyxminiproblems-problemphyxmini0615-lengthreadout}
  \lean{PhyXMiniProblems.ProblemPhyXMini0615.lengthReadout}
  \uses{def:physics:phyx-mini-0615:phyxminiproblems-problemphyxmini0615-lengthquantity}
  Read a physical length in a selected length unit
\end{definition}
\begin{proof}
  This is the declared model definition of length Readout.
\end{proof}
\begin{definition}[length In Meters]
  \label{def:physics:phyx-mini-0615:phyxminiproblems-problemphyxmini0615-lengthinmeters}
  \lean{PhyXMiniProblems.ProblemPhyXMini0615.lengthInMeters}
  \uses{def:physics:phyx-mini-0615:phyxminiproblems-problemphyxmini0615-lengthquantity, def:physics:phyx-mini-0615:phyxminiproblems-problemphyxmini0615-lengthreadout}
  Metre readout of a physical length
\end{definition}
\begin{proof}
  This is the declared model definition of length In Meters.
\end{proof}
\begin{definition}[length In Femtometers]
  \label{def:physics:phyx-mini-0615:phyxminiproblems-problemphyxmini0615-lengthinfemtometers}
  \lean{PhyXMiniProblems.ProblemPhyXMini0615.lengthInFemtometers}
  \uses{def:physics:phyx-mini-0615:phyxminiproblems-problemphyxmini0615-lengthquantity, def:physics:phyx-mini-0615:phyxminiproblems-problemphyxmini0615-lengthreadout}
  Femtometre readout of a physical length
\end{definition}
\begin{proof}
  This is the declared model definition of length In Femtometers.
\end{proof}
\begin{definition}[mass In Kilograms]
  \label{def:physics:phyx-mini-0615:phyxminiproblems-problemphyxmini0615-massinkilograms}
  \lean{PhyXMiniProblems.ProblemPhyXMini0615.massInKilograms}
  \uses{def:physics:phyx-mini-0615:phyxminiproblems-problemphyxmini0615-massquantity}
  Kilogram readout of a physical mass
\end{definition}
\begin{proof}
  This is the declared model definition of mass In Kilograms.
\end{proof}
\begin{definition}[energy In Joules]
  \label{def:physics:phyx-mini-0615:phyxminiproblems-problemphyxmini0615-energyinjoules}
  \lean{PhyXMiniProblems.ProblemPhyXMini0615.energyInJoules}
  \uses{def:physics:phyx-mini-0615:phyxminiproblems-problemphyxmini0615-energyquantity}
  Joule readout of a physical energy
\end{definition}
\begin{proof}
  This is the declared model definition of energy In Joules.
\end{proof}
\begin{definition}[energy In Electron Volts]
  \label{def:physics:phyx-mini-0615:phyxminiproblems-problemphyxmini0615-energyinelectronvolts}
  \lean{PhyXMiniProblems.ProblemPhyXMini0615.energyInElectronVolts}
  \uses{def:physics:phyx-mini-0615:phyxminiproblems-problemphyxmini0615-energyquantity}
  Electron-volt readout of a physical energy
\end{definition}
\begin{proof}
  This is the declared model definition of energy In Electron Volts.
\end{proof}
\begin{definition}[energy In Mega Electron Volts]
  \label{def:physics:phyx-mini-0615:phyxminiproblems-problemphyxmini0615-energyinmegaelectronvolts}
  \lean{PhyXMiniProblems.ProblemPhyXMini0615.energyInMegaElectronVolts}
  \uses{def:physics:phyx-mini-0615:phyxminiproblems-problemphyxmini0615-energyquantity, def:physics:phyx-mini-0615:phyxminiproblems-problemphyxmini0615-energyinelectronvolts}
  Mega-electron-volt readout of a physical energy
\end{definition}
\begin{proof}
  This is the declared model definition of energy In Mega Electron Volts.
\end{proof}
\begin{definition}[action In Joule Seconds]
  \label{def:physics:phyx-mini-0615:phyxminiproblems-problemphyxmini0615-actioninjouleseconds}
  \lean{PhyXMiniProblems.ProblemPhyXMini0615.actionInJouleSeconds}
  \uses{def:physics:phyx-mini-0615:phyxminiproblems-problemphyxmini0615-actionquantity}
  Joule-second readout of a physical action
\end{definition}
\begin{proof}
  This is the declared model definition of action In Joule Seconds.
\end{proof}
\begin{definition}[wave Number In Inverse Meters]
  \label{def:physics:phyx-mini-0615:phyxminiproblems-problemphyxmini0615-wavenumberininversemeters}
  \lean{PhyXMiniProblems.ProblemPhyXMini0615.waveNumberInInverseMeters}
  \uses{def:physics:phyx-mini-0615:phyxminiproblems-problemphyxmini0615-wavenumberquantity}
  Inverse-metre readout of a physical wave number
\end{definition}
\begin{proof}
  This is the declared model definition of wave Number In Inverse Meters.
\end{proof}
\begin{definition}[Particle Species]
  \label{def:physics:phyx-mini-0615:phyxminiproblems-problemphyxmini0615-particlespecies}
  \lean{PhyXMiniProblems.ProblemPhyXMini0615.ParticleSpecies}
  Particle species used in the simple nuclear model
\end{definition}
\begin{proof}
  This is the declared model definition of Particle Species.
\end{proof}
\begin{definition}[Nuclear Potential Model]
  \label{def:physics:phyx-mini-0615:phyxminiproblems-problemphyxmini0615-nuclearpotentialmodel}
  \lean{PhyXMiniProblems.ProblemPhyXMini0615.NuclearPotentialModel}
  Qualitative model selected for the nuclear potential
\end{definition}
\begin{proof}
  This is the declared model definition of Nuclear Potential Model.
\end{proof}
\begin{definition}[Plot Axis]
  \label{def:physics:phyx-mini-0615:phyxminiproblems-problemphyxmini0615-plotaxis}
  \lean{PhyXMiniProblems.ProblemPhyXMini0615.PlotAxis}
  The two axes drawn in the supplied barrier graph
\end{definition}
\begin{proof}
  This is the declared model definition of Plot Axis.
\end{proof}
\begin{definition}[Plot Axis Quantity]
  \label{def:physics:phyx-mini-0615:phyxminiproblems-problemphyxmini0615-plotaxisquantity}
  \lean{PhyXMiniProblems.ProblemPhyXMini0615.PlotAxisQuantity}
  Physical quantity assigned to an axis of the supplied graph
\end{definition}
\begin{proof}
  This is the declared model definition of Plot Axis Quantity.
\end{proof}
\begin{definition}[Potential Curve Shape]
  \label{def:physics:phyx-mini-0615:phyxminiproblems-problemphyxmini0615-potentialcurveshape}
  \lean{PhyXMiniProblems.ProblemPhyXMini0615.PotentialCurveShape}
  Shape of the potential-energy curve in the supplied graph
\end{definition}
\begin{proof}
  This is the declared model definition of Potential Curve Shape.
\end{proof}
\begin{definition}[Figure Annotation]
  \label{def:physics:phyx-mini-0615:phyxminiproblems-problemphyxmini0615-figureannotation}
  \lean{PhyXMiniProblems.ProblemPhyXMini0615.FigureAnnotation}
  Literal symbolic or numerical annotations visible in the primary image
\end{definition}
\begin{proof}
  This is the declared model definition of Figure Annotation.
\end{proof}
\begin{definition}[Rectangular Barrier Figure]
  \label{def:physics:phyx-mini-0615:phyxminiproblems-problemphyxmini0615-rectangularbarrierfigure}
  \lean{PhyXMiniProblems.ProblemPhyXMini0615.RectangularBarrierFigure}
  \uses{def:physics:phyx-mini-0615:phyxminiproblems-problemphyxmini0615-lengthquantity, def:physics:phyx-mini-0615:phyxminiproblems-problemphyxmini0615-energyquantity, def:physics:phyx-mini-0615:phyxminiproblems-problemphyxmini0615-plotaxis, def:physics:phyx-mini-0615:phyxminiproblems-problemphyxmini0615-plotaxisquantity, def:physics:phyx-mini-0615:phyxminiproblems-problemphyxmini0615-potentialcurveshape, def:physics:phyx-mini-0615:phyxminiproblems-problemphyxmini0615-figureannotation}
  The figure's energy-gap annotation is an independent field because its printed 1.0 MeV conflicts with the problem prose's physical 10.0 MeV deficit
\end{definition}
\begin{proof}
  This is the declared model definition of Rectangular Barrier Figure.
\end{proof}
\begin{definition}[Alpha Rectangular Barrier Setup]
  \label{def:physics:phyx-mini-0615:phyxminiproblems-problemphyxmini0615-alpharectangularbarriersetup}
  \lean{PhyXMiniProblems.ProblemPhyXMini0615.AlphaRectangularBarrierSetup}
  \uses{def:physics:phyx-mini-0615:phyxminiproblems-problemphyxmini0615-lengthquantity, def:physics:phyx-mini-0615:phyxminiproblems-problemphyxmini0615-massquantity, def:physics:phyx-mini-0615:phyxminiproblems-problemphyxmini0615-energyquantity, def:physics:phyx-mini-0615:phyxminiproblems-problemphyxmini0615-actionquantity, def:physics:phyx-mini-0615:phyxminiproblems-problemphyxmini0615-wavenumberquantity, def:physics:phyx-mini-0615:phyxminiproblems-problemphyxmini0615-particlespecies, def:physics:phyx-mini-0615:phyxminiproblems-problemphyxmini0615-nuclearpotentialmodel, def:physics:phyx-mini-0615:phyxminiproblems-problemphyxmini0615-rectangularbarrierfigure}
  Independent physical quantities in the alpha-tunneling model. In particular, the attenuation wave number and tunneling probability are independent fields; the governing laws below constrain them without defining either as the recorded answer
\end{definition}
\begin{proof}
  This is the declared model definition of Alpha Rectangular Barrier Setup.
\end{proof}
\begin{definition}[Matches Alpha Nucleus Scenario]
  \label{def:physics:phyx-mini-0615:phyxminiproblems-problemphyxmini0615-matchesalphanucleusscenario}
  \lean{PhyXMiniProblems.ProblemPhyXMini0615.MatchesAlphaNucleusScenario}
  \uses{def:physics:phyx-mini-0615:phyxminiproblems-problemphyxmini0615-alpharectangularbarriersetup}
  Qualitative physical scenario stated in the problem prose
\end{definition}
\begin{proof}
  This is the declared model definition of Matches Alpha Nucleus Scenario.
\end{proof}
\begin{definition}[Matches Supplied Barrier Figure]
  \label{def:physics:phyx-mini-0615:phyxminiproblems-problemphyxmini0615-matchessuppliedbarrierfigure}
  \lean{PhyXMiniProblems.ProblemPhyXMini0615.MatchesSuppliedBarrierFigure}
  \uses{def:physics:phyx-mini-0615:phyxminiproblems-problemphyxmini0615-energyinmegaelectronvolts, def:physics:phyx-mini-0615:phyxminiproblems-problemphyxmini0615-alpharectangularbarriersetup}
  Axes, labels, rectangular geometry, and energy levels read from the primary image. The printed 1.0 MeV annotation is deliberately not equated to the prose deficit, since doing so would make the supplied data inconsistent. This premise contains no tunneling probability or answer choice
\end{definition}
\begin{proof}
  This is the declared model definition of Matches Supplied Barrier Figure.
\end{proof}
\begin{definition}[Matches Problem Readouts]
  \label{def:physics:phyx-mini-0615:phyxminiproblems-problemphyxmini0615-matchesproblemreadouts}
  \lean{PhyXMiniProblems.ProblemPhyXMini0615.MatchesProblemReadouts}
  \uses{def:physics:phyx-mini-0615:phyxminiproblems-problemphyxmini0615-lengthinfemtometers, def:physics:phyx-mini-0615:phyxminiproblems-problemphyxmini0615-massinkilograms, def:physics:phyx-mini-0615:phyxminiproblems-problemphyxmini0615-energyinmegaelectronvolts, def:physics:phyx-mini-0615:phyxminiproblems-problemphyxmini0615-alpharectangularbarriersetup}
  Numerical physical data supplied by the problem prose. The final field states that the incident level is the specified deficit below the barrier top; it does not state a transmission probability
\end{definition}
\begin{proof}
  This is the declared model definition of Matches Problem Readouts.
\end{proof}
\begin{definition}[Uses Standard Reduced Planck Constant]
  \label{def:physics:phyx-mini-0615:phyxminiproblems-problemphyxmini0615-usesstandardreducedplanckconstant}
  \lean{PhyXMiniProblems.ProblemPhyXMini0615.UsesStandardReducedPlanckConstant}
  \uses{def:physics:phyx-mini-0615:phyxminiproblems-problemphyxmini0615-actioninjouleseconds, def:physics:phyx-mini-0615:phyxminiproblems-problemphyxmini0615-alpharectangularbarriersetup}
  The standard reduced Planck constant used by the numerical model
\end{definition}
\begin{proof}
  This is the declared model definition of Uses Standard Reduced Planck Constant.
\end{proof}
\begin{definition}[Matches Physlib Rectangular Barrier]
  \label{def:physics:phyx-mini-0615:phyxminiproblems-problemphyxmini0615-matchesphyslibrectangularbarrier}
  \lean{PhyXMiniProblems.ProblemPhyXMini0615.MatchesPhyslibRectangularBarrier}
  \uses{def:physics:phyx-mini-0615:phyxminiproblems-problemphyxmini0615-lengthinmeters, def:physics:phyx-mini-0615:phyxminiproblems-problemphyxmini0615-massinkilograms, def:physics:phyx-mini-0615:phyxminiproblems-problemphyxmini0615-energyinjoules, def:physics:phyx-mini-0615:phyxminiproblems-problemphyxmini0615-alpharectangularbarriersetup}
  Calibrate Physlib's scalar rectangular-barrier parameters by coherent-SI readouts of the dimensionful physical quantities. The absolute radial location of the two faces is left free because the image specifies only their separation
\end{definition}
\begin{proof}
  This is the declared model definition of Matches Physlib Rectangular Barrier.
\end{proof}
\begin{definition}[Has Physical Tunneling Parameters]
  \label{def:physics:phyx-mini-0615:phyxminiproblems-problemphyxmini0615-hasphysicaltunnelingparameters}
  \lean{PhyXMiniProblems.ProblemPhyXMini0615.HasPhysicalTunnelingParameters}
  \uses{def:physics:phyx-mini-0615:phyxminiproblems-problemphyxmini0615-lengthinmeters, def:physics:phyx-mini-0615:phyxminiproblems-problemphyxmini0615-massinkilograms, def:physics:phyx-mini-0615:phyxminiproblems-problemphyxmini0615-energyinjoules, def:physics:phyx-mini-0615:phyxminiproblems-problemphyxmini0615-actioninjouleseconds, def:physics:phyx-mini-0615:phyxminiproblems-problemphyxmini0615-wavenumberininversemeters, def:physics:phyx-mini-0615:phyxminiproblems-problemphyxmini0615-alpharectangularbarriersetup}
  Positivity and range conditions selecting a physical tunneling setup
\end{definition}
\begin{proof}
  This is the declared model definition of Has Physical Tunneling Parameters.
\end{proof}
\begin{definition}[Obeys Textbook Rectangular Barrier Tunneling]
  \label{def:physics:phyx-mini-0615:phyxminiproblems-problemphyxmini0615-obeystextbookrectangularbarriertunneling}
  \lean{PhyXMiniProblems.ProblemPhyXMini0615.ObeysTextbookRectangularBarrierTunneling}
  \uses{def:physics:phyx-mini-0615:phyxminiproblems-problemphyxmini0615-lengthinmeters, def:physics:phyx-mini-0615:phyxminiproblems-problemphyxmini0615-massinkilograms, def:physics:phyx-mini-0615:phyxminiproblems-problemphyxmini0615-energyinjoules, def:physics:phyx-mini-0615:phyxminiproblems-problemphyxmini0615-actioninjouleseconds, def:physics:phyx-mini-0615:phyxminiproblems-problemphyxmini0615-wavenumberininversemeters, def:physics:phyx-mini-0615:phyxminiproblems-problemphyxmini0615-alpharectangularbarriersetup}
  The two exact textbook relations used for the recorded multiple-choice computation in the one-dimensional, nonrelativistic, stationary rectangular- barrier model. In the classically forbidden region, ℏ² κ² = 2 m (U₀ - E). Exact matching of the stationary wavefunction and its derivative at both barrier faces gives T = (1 + U₀² sinh²(κ L) / (4 E (U₀ - E)))⁻¹ for 0 < E < U₀. Every term is an SI readout of a dimensionally typed quantity except the dimensionless ratios, exponent, and probability.
\end{definition}
\begin{proof}
  This is the declared model definition of Obeys Textbook Rectangular Barrier Tunneling.
\end{proof}
\begin{definition}[Answer Choice]
  \label{def:physics:phyx-mini-0615:phyxminiproblems-problemphyxmini0615-answerchoice}
  \lean{PhyXMiniProblems.ProblemPhyXMini0615.AnswerChoice}
  Labels of the four tunneling-probability choices
\end{definition}
\begin{proof}
  This is the declared model definition of Answer Choice.
\end{proof}
\begin{definition}[displayed Tunneling Probability]
  \label{def:physics:phyx-mini-0615:phyxminiproblems-problemphyxmini0615-displayedtunnelingprobability}
  \lean{PhyXMiniProblems.ProblemPhyXMini0615.displayedTunnelingProbability}
  \uses{def:physics:phyx-mini-0615:phyxminiproblems-problemphyxmini0615-answerchoice}
  Dimensionless probabilities printed beside the four choices
\end{definition}
\begin{proof}
  This is the declared model definition of displayed Tunneling Probability.
\end{proof}
\begin{definition}[recorded Dataset Answer]
  \label{def:physics:phyx-mini-0615:phyxminiproblems-problemphyxmini0615-recordeddatasetanswer}
  \lean{PhyXMiniProblems.ProblemPhyXMini0615.recordedDatasetAnswer}
  \uses{def:physics:phyx-mini-0615:phyxminiproblems-problemphyxmini0615-answerchoice}
  The answer label recorded by the source dataset
\end{definition}
\begin{proof}
  This is the declared model definition of recorded Dataset Answer.
\end{proof}
\begin{definition}[Rounds To Three Decimal Places]
  \label{def:physics:phyx-mini-0615:phyxminiproblems-problemphyxmini0615-roundstothreedecimalplaces}
  \lean{PhyXMiniProblems.ProblemPhyXMini0615.RoundsToThreeDecimalPlaces}
  value rounds to reported at three decimal places. The half-open interval implements ordinary nearest-thousandth rounding away from the upper tie
\end{definition}
\begin{proof}
  This is the declared model definition of Rounds To Three Decimal Places.
\end{proof}
\begin{definition}[Matches Displayed Tunneling Probability]
  \label{def:physics:phyx-mini-0615:phyxminiproblems-problemphyxmini0615-matchesdisplayedtunnelingprobability}
  \lean{PhyXMiniProblems.ProblemPhyXMini0615.MatchesDisplayedTunnelingProbability}
  \uses{def:physics:phyx-mini-0615:phyxminiproblems-problemphyxmini0615-alpharectangularbarriersetup, def:physics:phyx-mini-0615:phyxminiproblems-problemphyxmini0615-answerchoice, def:physics:phyx-mini-0615:phyxminiproblems-problemphyxmini0615-displayedtunnelingprobability, def:physics:phyx-mini-0615:phyxminiproblems-problemphyxmini0615-roundstothreedecimalplaces}
  A displayed choice agrees with the modeled probability to three decimals
\end{definition}
\begin{proof}
  This is the declared model definition of Matches Displayed Tunneling Probability.
\end{proof}
\begin{definition}[Is Unique Matching Tunneling Choice]
  \label{def:physics:phyx-mini-0615:phyxminiproblems-problemphyxmini0615-isuniquematchingtunnelingchoice}
  \lean{PhyXMiniProblems.ProblemPhyXMini0615.IsUniqueMatchingTunnelingChoice}
  \uses{def:physics:phyx-mini-0615:phyxminiproblems-problemphyxmini0615-alpharectangularbarriersetup, def:physics:phyx-mini-0615:phyxminiproblems-problemphyxmini0615-answerchoice, def:physics:phyx-mini-0615:phyxminiproblems-problemphyxmini0615-matchesdisplayedtunnelingprobability}
  A choice is the unique displayed value matching the modeled probability
\end{definition}
\begin{proof}
  This is the declared model definition of Is Unique Matching Tunneling Choice.
\end{proof}
\begin{lemma}[incident kinetic energy megaelectronvolts eq]
  \label{lem:physics:phyx-mini-0615:phyxminiproblems-problemphyxmini0615-incident-kinetic-energy-megaelectronvolts-eq}
  \lean{PhyXMiniProblems.ProblemPhyXMini0615.incident_kinetic_energy_megaelectronvolts_eq}
  \uses{def:physics:phyx-mini-0615:phyxminiproblems-problemphyxmini0615-energyinmegaelectronvolts, def:physics:phyx-mini-0615:phyxminiproblems-problemphyxmini0615-alpharectangularbarriersetup, def:physics:phyx-mini-0615:phyxminiproblems-problemphyxmini0615-matchesproblemreadouts}
  The prose energy separation implies an incident energy of 20 MeV
\end{lemma}
\begin{proof}
  The conclusion follows from the governing relations and figure readouts named in the statement of incident kinetic energy megaelectronvolts eq.
\end{proof}
\begin{lemma}[tunneling probability rounds to 0 014]
  \label{lem:physics:phyx-mini-0615:phyxminiproblems-problemphyxmini0615-tunneling-probability-rounds-to-0-014}
  \lean{PhyXMiniProblems.ProblemPhyXMini0615.tunneling_probability_rounds_to_0_014}
  \uses{def:physics:phyx-mini-0615:phyxminiproblems-problemphyxmini0615-alpharectangularbarriersetup, def:physics:phyx-mini-0615:phyxminiproblems-problemphyxmini0615-matchesproblemreadouts, def:physics:phyx-mini-0615:phyxminiproblems-problemphyxmini0615-usesstandardreducedplanckconstant, def:physics:phyx-mini-0615:phyxminiproblems-problemphyxmini0615-hasphysicaltunnelingparameters, def:physics:phyx-mini-0615:phyxminiproblems-problemphyxmini0615-obeystextbookrectangularbarriertunneling, def:physics:phyx-mini-0615:phyxminiproblems-problemphyxmini0615-roundstothreedecimalplaces}
  The dimensionful data and the exact rectangular-barrier tunneling laws put the transmission probability in the interval that rounds to 0.014
\end{lemma}
\begin{proof}
  The conclusion follows from the governing relations and figure readouts named in the statement of tunneling probability rounds to 0 014.
\end{proof}
% --- Archon declaration topology end ---
% --- Archon physics formalization source end ---
